\chapter{Conclusion}

This thesis set out to determine whether bounded Android permission evidence could support a useful privacy-review product without being misrepresented as automated legal judgement. The conclusion answers the four research questions, integrates the architectural and methodological contributions, states the final product position, and identifies the next evidence required. Operational deployment details already analysed in Chapter 6 are not repeated here.

\section{Answers to the Research Questions}

\textbf{RQ1: deterministic and safe transformation.} Permission summaries are accepted through a conservative evidence contract, rejected when identifiers, numbers, sources, windows, timestamps, or contexts are unsafe, and evaluated through a deterministic rule-pack engine. Temporal isolation addresses replay and overlap within the limits of aggregate windows.

The answer rests on the recorded suites rather than on nominal examples: runtime validation treats TypeScript declarations as developer guidance rather than as a security boundary, so a rejected observation does not become a normal finding, and determinism extends to state, where an identical record does not accumulate, the retention watermark cannot move backwards, and cross-window evaluation admits only stored windows that end no later than the candidate window's start. Two qualifications bound the answer. The watermark belongs to in-memory engine state and is not currently witnessed by a dedicated test, and the admission test prevents the tested forms of aggregate amplification without guaranteeing that retained windows are mutually disjoint. Nothing here reconstructs information that the producer did not supply.

\textbf{RQ2: bounded communication and presentation.} Each finding preserves provenance, temporal scope, rule-pack identity, rationale, and missing-context questions. The interface distinguishes needs-review, evidence-gap, and no-technical-concern states, and repeatedly states that technical output is not a legal verdict. Synthetic demonstrations remain labelled throughout the data path.

Bounded communication is implemented as an information architecture rather than as a disclaimer appended to an absolute score, and the recorded device render shows that the path exists on the tested handset. It does not show that a population understands the distinction between technical concern and legal compliance, which would require a task-based participant study. RQ2 is therefore answered at the level of traceable interface design and bounded device evidence, not validated human comprehension.

\textbf{RQ3: regional extensibility.} Jurisdiction-dependent material is loaded from registered, versioned packs behind a stable contract. The EU GDPR pack is the evaluated legal objective; the Research Baseline demonstrates software switching without claiming another jurisdiction. New regional packs require their own legal and empirical validation.

Two forms of extensibility are separated here. Software extensibility means that another admitted pack can supply identifiers, labels, thresholds, references, rationales and questions without rewriting navigation, persistence or Android integration; substantive legal extensibility means that those contents are correct for a jurisdiction and use context. The project demonstrates the first and defines governance requirements for the second, and a finding stores the pack and rule version that produced it so that a later update need not erase historical meaning.

\textbf{RQ4: evidence strength and external validity.} Deterministic, adversarial, and curated mutation suites support conformance and sensitivity to the registered semantic weakenings. The hash-pinned \FDroidCensusN{}-APK census supports processing of heterogeneous open-source artifacts. The complete \DevicePackagesN{}-package device denominator, including \FailedDeviceWorkflowsN{} retained failures and \PssObservedN{} PSS observations, supports only bounded single-device engineering statements. The FlowDroid receipt supports interoperability analysis while preserving absent XML as missing evidence. None of these results establishes Android-population prevalence, legal compliance, human comprehension, unrestricted Android visibility, long-run reliability, production signing, or store acceptance.

The answer is a claim map rather than a maturity slogan: logical suites support engine conformance to specified rules, package inspection supports statements about identifiers, SDK targets, permissions, backup configuration and incorporated resources, physical checks support only the recorded handset, build and interaction interval, and visual review supports an expert judgement about coherence for controlled submission work. This structure makes negative results usable, because a missed observation, a misunderstood label or a revised legal mapping narrows or changes one claim without invalidating the others.

\section{Contributions and Research Significance}

The three contributions defined in Chapter 1 and evaluated in Chapter 5 are restated here with the strength that the recorded evidence supports.

\textbf{C1: cross-layer evidence-bounded contracts.} The governing principle is that every conclusion must retain the source, temporal scope, rule identity, and known missing context needed to interpret it, and that the permitted output domain excludes legal verdicts. This principle is not a wrapper around the implementation. It shapes data types, admission validation, storage, interface language, testing, and release documentation, and it is what prevents individually working features from combining into a claim the system cannot support. Its concrete form is a closed identifier and source vocabulary, safe-integer and window validation, replay semantics, structured evidence gaps, and a typed output space that admits review obligations but not infringement or compliance findings. The pipeline is deliberately deterministic so that a reviewer can reconstruct why a record was accepted and why a finding changed. The regulation-pack contract belongs to this contribution as well: it locates jurisdictional interpretation in a reviewable, versioned component, which is why a finding stores the pack and rule identity that produced it. The contribution is the architecture and the governance model, not a claim that arbitrary legal systems can be translated automatically.

\textbf{C2: bounded temporal combination and state continuity.} Temporal claims are made at a stated abstraction boundary and are supported only inside it. Exact-boundary cases, all 720 permutations of a six-event fixture, replay and overlap constructions, pack-switch rejection, bounded-ledger behaviour, malformed-input rejection, and 10{,}000 export-and-restore cases show order-independent and recovery-preserving behaviour for the recorded conditions. The contribution is the explicit statement of those conditions and of the conservative treatment applied when they do not hold. Where a state object is only partially exported, or where aggregate evidence cannot identify the events it summarises, the thesis records the loss rather than extending the recovery claim. What the design adds is therefore a bounded and inspectable continuity semantics, not the general novelty of replay suppression or bounded storage.

\textbf{C3: layered, diagnostic evaluation.} The evaluation record connects artefacts to claims without collapsing unlike evidence into one score. Deterministic fixtures, nine registered safety properties with curated mutants, typed-policy fixtures, package and device receipts, missingness summaries, a static-analysis record with no emitted artifact, and unresolved review gates are reported with separate denominators and separate conclusions. The most transferable finding is diagnostic rather than quantitative: executing a test, distinguishing a wrong implementation, and establishing external validity are three different achievements, and the gap between them is where verification claims most commonly overreach. This contribution is what keeps the evaluation from inflating its own evidence, and its value is methodological: because the record reports each evidence object with its own denominator, names the tier of the oracle behind each suite, and retains failures and missing values instead of absorbing them, a conformance result cannot quietly become an accuracy claim and a curated mutant set cannot quietly become a mutation score. Retaining logs, recording hashes and maintaining a claim matrix are necessary hygiene for that discipline; the contribution lies in the separation they make auditable rather than in the record-keeping itself.

Together, these contributions support a form of trustworthy incompleteness. Privacy Lens does not solve every acquisition, legal, or usability problem, but it makes those unsolved parts structurally difficult to hide. In an accountability tool, that property is itself useful. A system that clearly refuses an unsupported verdict can better support responsible human review than a more ambitious classifier whose authority is unclear. The Android application, the three-destination interface, local storage, and the reproducibility toolchain are the carriers of these contributions rather than separate contributions of their own.

\section{Final Product Position}

Privacy Lens \PrivacyLensVersion{} is the current research prototype release identity. Its bounded engineering evidence supports controlled review and reproduction; it does not establish user comprehension, legal validity, production readiness, or broad device reliability. The release artefacts use package \texttt{com.zhihengzhang.privacylens}, target SDK 36, and request no sensitive runtime permission. They are accompanied by a privacy-policy draft, user guide, product audit, store-readiness checklist, checksums, and physical-device screenshots.

The candidate is not yet a publicly publishable production release. Production signing, stable public policy and support endpoints, console declarations, internal-track evaluation, and store review remain external owner actions. This distinction is part of the result rather than an administrative footnote: trustworthy privacy software must disclose where its evidence ends.

\section{Implications for Privacy-Engineering Practice}

The first lesson is that provenance must survive every transformation. Capturing a source label at ingestion is insufficient if storage, exports, or screenshots omit it. Provenance should be treated as part of the value being computed. Tests should therefore follow it from observation through decision and presentation.

The second lesson is that missing context requires a first-class state. Many privacy questions cannot be answered from technical frequency. Returning a neutral or positive result when purpose, authority, or acquisition coverage is unknown creates hidden assumptions. An evidence-gap state makes the limitation actionable: it can ask for a privacy notice, controller record, consent basis, or authorised trace rather than fabricating certainty.

The third lesson is that extensibility needs governance. An interface or schema that can load new legal text is technically flexible but may accelerate error if sources, versions, review, and updates are uncontrolled. The pack boundary is valuable precisely because it creates a unit that can be admitted, tested, rejected, and deprecated.

The fourth lesson is that polished design should reduce ambiguity rather than conceal it. Visual hierarchy can foreground evidence strength, state, and next action. Calm colour and concise language can lower cognitive burden. Yet critical caveats must remain visible. The ethical objective is not to maximise engagement or confidence; it is to help the user reach an appropriately bounded understanding.

The fifth lesson is that store readiness is an evidence package. An AAB is necessary but not sufficient. Signing ownership, manifest reconciliation, data declarations, policy endpoints, screenshots, accessibility, update testing, and support processes all affect whether the distributed product matches the evaluated one. Treating these as afterthoughts would break the thesis's claim chain at the point of deployment.

The final lesson is that disciplined limitation can increase, rather than reduce, practical value. A reviewer can act on a clearly sourced question, ask for missing evidence, and revisit the decision after a pack update. An unsupported compliance badge offers none of those affordances. Privacy Lens is most credible when it remains specific about what it observed, how it reasoned, and what a human must still decide.

\section{Future Work}

Future work is separated by the kind of evidence or operational authority it requires. These activities should retain the thesis's claim-to-evidence discipline. Additional features are valuable only when their authority, provenance, and validation remain visible.

\subsection{Immediate Validation Priorities}

The immediate priorities are independent validation of the GDPR pack, enabled-service accessibility testing including TalkBack and large text, and task-based participant evaluation of comprehension and reliance. A multi-OEM replication should repeat the acquisition and interaction protocol on supported Android versions and retain missing observations rather than converting them to negatives. Production signing, stable policy and support endpoints, Play Console declarations, internal-track delivery, and store review remain owner-controlled release gates.

\subsection{Research Extensions}

The programme should extend the existing formal and property-based techniques for evidence admission and replay. The registered suite covers idempotence, bounded storage, source isolation, reordered independent windows, pack isolation, and missing outputs. Monotonic retention of the aggregation watermark is a design invariant with source-level support but no dedicated witness, and it should be added to the suite rather than listed as already covered; the same applies to the mutually-overlapping prior-window case that the current admission test does not exclude. The curated mutants should be complemented by automatically generated and independently designed mutants. A separately implemented oracle would reduce the correlated-defect risk of deriving implementation and expected values from one specification.

Interface research should compare alternative presentations of uncertainty. For example, a three-state label, an evidence-strength scale, and a question-led presentation may produce different reliance patterns. The comparison should measure comprehension and decisions, not only preference. Accessibility evaluation should be incorporated from recruitment through reporting rather than added as a final conformance scan.

The pack framework could support signed declarative updates after its threat model is tested. This would require publisher identity, schema versioning, offline recovery, rollback protection, and clear user communication. Remote executable rules would introduce unnecessary risk; the safest extension preserves a narrow validated data contract.

Finally, future scholarship should study the relationship between accountability prompts and organisational action. A technically clear question has value only if users or reviewers can obtain the missing policy, consent, retention, or controller evidence. Research with data-protection practitioners could examine which prompts accelerate review, which create noise, and how findings should integrate with established assessment records without becoming an unreviewed legal decision.

\subsection{Productisation}

Productisation should improve authorised event-level acquisition without implying unrestricted cross-application visibility, exercise long-duration scheduling and recovery across Android implementations, and establish an operational vulnerability, support, retention, and update process. The pack framework may support a genuinely reviewed regional pack only after versioned authority, independent domain review, localisation, rollback, and deprecation procedures are in place. These are staged gates, not capabilities attributed to the submitted prototype.

\section{Closing Statement}

The central finding is that privacy accountability on Android is strengthened not by presenting the strongest possible conclusion, but by presenting the strongest conclusion the available evidence actually supports. Privacy Lens operationalises that principle in code, interface, evaluation, and thesis structure. Appendix~\ref{app:claims} records each core claim with its evidence type and its boundary, so that the same discipline can be applied to the thesis itself.
